\section{Discussion}\label{sec:discussion}

\paragraph{Reporting tail risk honestly.} The coverage study has a direct implication for how return periods of catastrophic cyber losses should be stated. With records of a few decades, the standard delta-method interval and its parametric-bootstrap substitute are not merely too narrow; they are centred on an estimate that is typically a fraction of the truth, and their misses all fall on the side that understates the loss. Statements of the form ``a loss of at least $x$ is expected once every $T$ years'' should therefore be accompanied by profile-likelihood intervals, and prudential uses should focus on the upper limit. Three further practices follow from our results. First, return levels should be reported across a range of thresholds, because, as \cref{sec:single,sec:coverage} show, threshold choice can bias estimates by a factor of two even with long records, and no interval method repairs that. Second, because the width of an honest interval shrinks rapidly with the number of exceedances (\cref{fig:coverage}c), pooling incident data across institutions and jurisdictions in supervisory or consortium databases has a high statistical return. Third, decision-makers should be told plainly what thin data imply: with 15 to 20 exceedances, an honest 95\% interval for the 1-in-1,000-year loss spans two to three orders of magnitude. A precise-looking estimate from such data should be read as a symptom of an overconfident method, not as information about the tail. Since capital and insurance premiums calibrated to a median-biased estimate will typically be too low, a supervisory margin over point estimates from short records is justified on statistical grounds alone.

\paragraph{Finality as a policy lever.} The attacker model suggests that the banking core is protected less by the difficulty of breaching it than by the reversibility of what an attacker can extract. Investments in payment recall, cross-institution freezing, anomaly-triggered settlement delays and international cooperation on tracing therefore carry a deterrence value beyond the losses they recover directly, because they remove the venue from the attacker's choice set (\cref{cor:breakeven}). Two corollaries deserve attention. First, the growth of instant, irrevocable payment systems raises finality inside the banking system itself. Our model predicts that, unless recovery mechanisms such as payee verification, reimbursement schemes or conditional delays restore the recovery probability, attacker effort will migrate towards these rails; this is a prediction of the model, not an empirical finding. Second, the asymmetry across venues is large. Because expected punishment and attack costs are small relative to retained value in crypto-asset venues, recovery alone deters attacks on the exchange only above a recovery probability of \breakEvenExchange{} (from \exchangeKappa{} in our calibration) and on the bridge only above \breakEvenBridge{} (from \bridgeKappa{}), against \breakEven{} for the bank. Issuer freeze functions, exchange cooperation and circuit breakers on bridges raise recovery, but at these levels attribution and sanction, the other terms of \cref{eq:payoff}, must do much of the work.

\paragraph{Interconnection and third-party concentration.} The contagion results imply that systemic cyber risk is strongly non-linear in interbank exposure. Below the critical exposure the system absorbs even large shocks with modest amplification; slightly above it, the same shock can destroy most of the system's equity. This argues for monitoring the spectral radius of the exposure matrix, or tractable proxies of it, as a macroprudential indicator, because the distance to criticality is more informative than the level of exposure itself. The ranking results argue for network-based measures of systemic importance alongside size-based ones, since a mid-sized periphery bank can transmit more loss per dollar of damage than most core banks. Finally, the shared-vendor scenario shows why third-party concentration deserves the supervisory attention it now receives \citep{dora2022}: in subcritical conditions a vendor compromise is benign, which is precisely why its danger is easy to underestimate, while near criticality it is the scenario that brings the system down. Stress tests that combine operational concentration with balance-sheet exposure, rather than treating them separately, capture this interaction.

\section{Limitations and extensions}\label{sec:limits}

\paragraph{Calibration.} Every parameter in this paper is illustrative. A calibrated version would estimate frequency and severity from operational-loss and incident databases, crypto-venue parameters from theft and recovery statistics such as those compiled by blockchain analytics firms, and the network from supervisory exposure data. The methodological conclusions depend on the tail index, the number of exceedances and the network spectrum rather than on dollar values, but their quantitative force in any application must be re-established with that application's data.

\paragraph{Stationarity and independence.} We treat years as independent and parameters as constant. Real cyber risk trends upward and clusters in campaigns that exploit common vulnerabilities. Both features would widen honest intervals further, so our coverage results are, if anything, optimistic. Self-exciting frequency models and time-varying severity are natural extensions.

\paragraph{Residual undercoverage.} With about 15 exceedances, profile-likelihood intervals still cover only about 90\%. A Bartlett correction of the likelihood-ratio statistic \citep{bartlett1937} or a bootstrap calibration of its distribution could close the gap; both are straightforward extensions of our code.

\paragraph{Threshold choice.} We fix thresholds by percentile rules, as practitioners often do. Automatic threshold selection, or models that span body and tail jointly, may reduce the bias documented in \cref{sec:single}.

\paragraph{Contagion mechanics.} Linear DebtRank passes equity losses through to claim values in full, which makes it an upper-bound stress test. Clearing models in the tradition of \citet{eisenberg2001} transmit losses only on default and imply much smaller amplification \citep{glasserman2015}; reality likely lies in between. We also omit liquidity and run channels, which \citet{duffie2019} and \citet{eisenbach2022} show can be first-order for cyber events that disrupt payments.

\paragraph{Attacker objectives.} The attacker model is static and risk-neutral, and it applies to financially motivated actors. State-sponsored attackers who seek disruption rather than profit are not deterred by finality, and they are arguably the more relevant threat for systemic events. The finality argument explains where profit-seeking attackers go; it does not imply that the banking core is safe from attackers with other objectives.

\paragraph{The compound-sum correction.} \Cref{app:shift} is derived heuristically and validated numerically; a proof under general regularity conditions, with a rate for the remainder, is left for future work.

\section{Conclusion}\label{sec:conclusion}

How bad, and how often? For the catastrophic cyber losses that worry supervisors, the honest answer from realistic data is a wide interval, and the intervals in common use are not honest. When a 1-in-1,000-year loss is estimated from 30 years of data, standard intervals miss the truth roughly 40\% of the time, always on the side that understates the loss, because the estimate they are centred on is typically a fraction of the true value. Profile-likelihood intervals restore most of the lost coverage and score better under a proper scoring rule, at the cost of upper limits that can be orders of magnitude above lower ones; that cost is the information content of the data, not a flaw of the method. Around this core result, we showed that the economics of attack are governed by finality, so that the banking core's main protection is the reversibility of what is stolen, and that losses which do land on banks propagate through interbank exposures with a sharp phase transition and with an importance ranking that size alone does not capture. All three findings point in the same direction: the quantities that matter most for systemic cyber risk are the ones a short record and a size-based view are least able to see.
